\section{BFieldExpansion: RK4 tracking}
\label{sec:bfield-expansion}
\label{chapter:bfield-expansion}

The \texttt{BFieldExpansion} element implements the field and vector-potential
construction in Secs.~\ref{sec:magnetic-field-reconstruction}
and \ref{sec:magnetic-vector-potential}, with polynomial longitudinal
profiles. It advances canonical coordinates using classical RK4, which is
not an exactly symplectic integrator.

\subsection{Polynomial profiles and additional strengths}

The current implementation selects straight geometry for \(h=0\) and
supports curved geometry for \(h>10^{-4}\,\mathrm{m}^{-1}\).
The quantities returned as \texttt{Bx}, \texttt{By}, \texttt{Bs},
\texttt{Ax}, \texttt{Ay}, \texttt{As}, and \texttt{phi} by
\texttt{get\_field} use the normalization in Eq.~\eqref{bfe:normalization}.

The polynomial coordinate runs from \texttt{s\_start} to
\(\texttt{s\_start}+L\), where \(L\) is the element length. It need not
coincide with the accumulated particle coordinate \texttt{s} in a line.
The field-evaluation argument \texttt{s\_local} is entrance-relative, so
that the profiles are evaluated at
\(s=\texttt{s\_start}+\texttt{s\_local}\).

With only polynomial input coefficients and unit scale,
\begin{equation}
 N_n(s)=\sum_{j=0}^{d_n}\texttt{knc}_{nj}s^j,\qquad
 S_n(s)=\sum_{j=0}^{d_s}\texttt{ksc}_{nj}s^j,\qquad
 K_s(s)=\sum_{j=0}^{d_\ell}\texttt{ksolc}_{j}s^j.
 \label{bfe:coefficients}
\end{equation}
The transverse factorial is included in Eq.~\eqref{bfe:profiles}; there
is no longitudinal factorial in Eq.~\eqref{bfe:coefficients}.
The suffix \texttt{c} denotes coefficients. The three polynomial degrees
\(d_n\), \(d_s\), and \(d_\ell\) are independent. Transverse rows of
different lengths are padded within their own matrix at construction;
omitted profiles remain empty and contribute zero. Every supplied
\texttt{ksolc} coefficient is used. Its scalar-potential seed is
\begin{equation}
 \phi_0^{\mathrm{sol}}(s)
 =-\int_0^s K_s(u)\,\dee u
 =-\sum_{j=0}^{d_\ell}\frac{\texttt{ksolc}_{j}}{j+1}s^{j+1}.
\end{equation}
The additional degree is reserved internally, so users do not pad the
longitudinal profile or the transverse matrices for this integral.
In particular, a constant longitudinal field needs only
\texttt{ksolc=[ks]}. The read-only array \texttt{ksoll} contains its
longitudinal integral over the element interval, including \texttt{kscale}.

Uniform scalar strengths and integrated strengths are incorporated before
constructing the potential. For example,
\begin{equation}
 N_n(s)=\texttt{kscale}\left[
 \sum_j\texttt{knc}_{nj}s^j+k_n+\frac{\texttt{knl}_n}{L}\right],
 \label{bfe:total-profile}
\end{equation}
with the analogous expression for the skew profile. The independent
\(k_n\) and skew scalar strengths exist through octupole order; missing
entries contribute zero. The factor \texttt{kscale} also multiplies
\(K_s\), all reconstructed fields, and both potentials. Nonzero integrated
strengths require \(L\ne0\).

\subsection{Expansion order and convergence}

The parameter \texttt{num\_phi} specifies the maximum retained power of
\(y\) in field evaluation. An additional scalar-potential coefficient is
stored to obtain \(b_y\) by differentiation; vector potentials are evaluated
through the requested order. For a nonterminating expansion, the
infinite-series identities above hold only to the orders retained by the
corresponding series. Finite truncation is an approximation to the
three-dimensional Maxwell field.

The default \texttt{num\_phi='auto'} retains the complete polynomial
reconstruction, including the vector potential, in straight geometry. Its
choice accounts for allocated polynomial degrees and transverse orders,
as well as the independent scalar strengths, including coefficients that
are initially zero. In curved geometry it adds two vertical orders, retaining
the full expansion through first order in \(h\). Higher powers of \(h\)
are generally incomplete: the curved recurrence need not terminate.
Increasing an explicit \texttt{num\_phi} and checking convergence over the
aperture of interest is therefore necessary when higher curvature terms matter.
The allocated shapes and truncation order remain fixed after construction,
including capacity for explicitly supplied zero coefficients. Evaluation skips
unused orders and degrees, with the populated bounds recomputed after input
updates.

\subsection{Sparsity of the coefficient cache}

Each application of the recurrence \eqref{bfe:recurrence} moves two rows
down in \(i\) and, in the curved basis of
Eq.~\eqref{bfe:curved-coefficients}, two powers to the left in \(m\). The
nonzero coefficients therefore form a diagonal band: a normal seed
populates the odd rows, a skew or longitudinal seed the even rows, and a
row of a given parity may be entirely empty. For a curved element with
only dipole, quadrupole, and sextupole strengths, the automatic order
allocates nine rows and thirteen powers of \(w\), of which seven cells are
nonzero.

The evaluator records the populated range of \(m\) for every row after
each coefficient update, skips empty rows, and visits only that range.
Each coefficient polynomial in \(s\) is evaluated once per field
evaluation: row \(i\) supplies \(\phi\), \(b_x\), \(b_s\), and the vector
potential at vertical order \(i\), and \(b_y\) at order \(i-1\) through
Eq.~\eqref{bfe:field}. The cost of a field evaluation thus scales with the
number of populated coefficients rather than with the allocated capacity,
and a lower explicit \texttt{num\_phi} only pays off through the rows it
removes from the band.

\subsection{Accuracy in curved geometry}

The seed rows are evaluated as polynomials in \(x\), as described in
Sec.~\ref{sec:curved-seed-rows}, so the midplane field \(b_y(x,0,s)\) is
exact to roundoff at every multipole order and for every admissible
\(h\). The off-midplane rows are evaluated in powers of \(w=1+hx\) and
retain a relative roundoff of order \(\epsilon\,(hx)^{-(n-3)}\) for a
normal and \(\epsilon\,(hx)^{-(n-2)}\) for a skew multipole of order
\(n\). The lower bound \(h>10^{-4}\,\mathrm{m}^{-1}\) selects the curved
code path; it is not an accuracy guarantee for high orders. Compare
\texttt{get\_field} with a straight element of the same strengths when
combining a small curvature with octupole or higher orders: the physical
difference scales with \(hx\), whereas roundoff does not.

There are three separate approximations to control: the input profile fit,
the transverse expansion order, and the integration step size. Higher
longitudinal derivatives enter Eq.~\eqref{bfe:recurrence}, so agreement of
an on-axis fit alone does not establish off-axis accuracy. Increasing the
number of tracking steps cannot restore fringe terms absent from the field
model. For piecewise profiles, derivatives and vector-potential matching at
the joins also need to be checked.

\subsection{Hamiltonian and equations of motion}

We use the canonical variables \(z=(x,p_x,y,p_y,\tau,p_\tau)\) defined in
Eqs.~\eqref{eq:coord1}--\eqref{eq:coord2}, with
\(\tau=\zeta/\beta_0=s/\beta_0-ct\), and the species factor
\(\chi=(q/q_0)(m_0/m)\). Let
\begin{equation}
 \begin{aligned}
 D&=1+\delta=\sqrt{1+2p_\tau/\beta_0+p_\tau^2}, &
 \pi_s&=\sqrt{D^2-\pi_x^2-\pi_y^2},\\
 \pi_x&=p_x-\chi a_x, & \pi_y&=p_y-\chi a_y.
 \end{aligned}
 \label{bfe:kinetic-momentum}
\end{equation}
Here \(\bm\pi\) is the normalized kinetic momentum. The forward branch
\(\pi_s>0\) is used, away from longitudinal turning points.
Specializing the accelerator Hamiltonian
\eqref{eq:hamiltonian_htau} to a static magnetic field gives
\begin{equation}
 H=\frac{p_\tau}{\beta_0}-w(\pi_s+\chi a_s).
 \label{bfe:hamiltonian}
\end{equation}
The Hamiltonian formalism and canonical normalization are discussed in
Ref.~\cite{wolski2014beam}. No paraxial expansion of the square root is
made here.

In the implemented gauge \(a_y=0\), Hamilton's equations are
\begin{align}
 \frac{\dee x}{\dee s}&=w\frac{\pi_x}{\pi_s}, &
 \frac{\dee y}{\dee s}&=w\frac{\pi_y}{\pi_s},\\
 \frac{\dee p_x}{\dee s}&=h(\pi_s+\chi a_s)
   +w\chi\left(\frac{\pi_x}{\pi_s}\partial_x a_x+\partial_x a_s\right), &
 \frac{\dee p_y}{\dee s}&=w\chi\left(
       \frac{\pi_x}{\pi_s}\partial_y a_x+\partial_y a_s\right),\\
 \frac{\dee\tau}{\dee s}&=\frac{1}{\beta_0}
                -w\frac{p_\tau+1/\beta_0}{\pi_s}, &
 \frac{\dee p_\tau}{\dee s}&=0.
 \label{bfe:equations}
\end{align}
Energy, and hence \(D\), is constant because there is no electric field.
The spatial Hamiltonian itself can depend explicitly on \(s\) and need
not be constant along the trajectory.

\subsection{Numerical integration}

Both geometries use classical RK4~\cite{HairerNorsettWanner1993}. For
\(\dee z/\dee s=f(s,z)\), a step of length
\(\Delta s=L/\texttt{num\_integration\_steps}\) is
\begin{align}
 k_1&=f(s_n,z_n),\\
 k_2&=f(s_n+\Delta s/2,z_n+\Delta s\,k_1/2),\\
 k_3&=f(s_n+\Delta s/2,z_n+\Delta s\,k_2/2),\\
 k_4&=f(s_n+\Delta s,z_n+\Delta s\,k_3),\\
 z_{n+1}&=z_n+\frac{\Delta s}{6}(k_1+2k_2+2k_3+k_4).
\end{align}
The potential and its derivatives are reevaluated at every stage. For a
fixed smooth field and a trajectory bounded away from \(\pi_s=0\), the
local integration error is \(\mathcal O(\Delta s^5)\), and the global
error over a fixed length is \(\mathcal O(\Delta s^4)\). This statement
concerns integration error, separately from field truncation and fitting.

Classical RK4 does not preserve the symplectic form exactly, even when
applied to Hamilton's equations. Nor does it generally preserve phase-space
volume exactly. The distinction between Hamiltonian flow and a symplectic
discrete map is discussed in Ref.~\cite{HairerLubichWanner2006}.
Here \(p_\tau\) remains constant because its right-hand side is identically
zero; this fact alone does not imply symplecticity. Backtracking reverses
\(\Delta s\) and the integration interval, but RK4 is not exactly reversible
at finite step size.

The element integrates the orbit only. It emits no synchrotron radiation
when radiation is enabled for the line, and it does not transport spin:
the spin vector passes through unchanged and no error is raised during
line tracking. Only the standalone \texttt{track} method of the element
rejects particles with a nonzero spin.

\subsection{Canonical and kinetic momenta at boundaries}

The flag \texttt{pkin\_const} selects how transverse momenta are transferred
between vector-potential conventions at the entrance and exit. It does not
change RK4 or its order of accuracy. Let
\(\alpha_{x,y}=\chi a_{x,y}\); the particle attributes \texttt{ax} and
\texttt{ay} store these species-scaled quantities. Then
\(\pi_{x,y}=p_{x,y}-\alpha_{x,y}\), as exposed by
\texttt{kin\_px} and \texttt{kin\_py}.

With \texttt{pkin\_const=False}, the default, the entrance canonical
momenta are used without conversion. At the exit, the element leaves the
momenta in its gauge and stores the exit \(\alpha_x,\alpha_y\).
Consequently, keeping \(p_x,p_y\) fixed while changing the transverse
vector potential at a join changes the kinetic momenta. To initialize a
particle with prescribed kinetic momenta inside a nonzero field, use
\begin{equation}
 p_x=\pi_x+\chi a_x(x,y,s_{\mathrm{in}}),\qquad
 p_y=\pi_y+\chi a_y(x,y,s_{\mathrm{in}}),
\end{equation}
and set the stored \texttt{ax}, \texttt{ay} consistently.

With \texttt{pkin\_const=True}, the entrance conversion is
\begin{equation}
 p_{x,y}^{\mathrm{in}}=p_{x,y}^{\mathrm{old}}
          +\chi a_{x,y}^{\mathrm{in}}-\alpha_{x,y}^{\mathrm{old}}.
\end{equation}
At the exit the element subtracts \(\chi a_{x,y}^{\mathrm{out}}\) and
sets the stored transverse vector potential to zero. These conversions
preserve kinetic momenta across the changes of potential convention.
They do not supply missing physical fringe fields or make the finite-step
map symplectic. For a smooth piecewise magnetic model, the field profiles,
their derivatives, and the choice of boundary convention must describe a
consistent physical join.
